%% In this file, the package imports are listed, style and formatting rules are declared, and in the last couple rows the contents.tex is included, thus this file is the entry point for rendering.

% General rules
% --------------------------------------------------

% twoside: printing margins are set two sided printing --> you can change it to "oneside" (but, please, protect environment and don't do that especially because the "Code of Studies and Exams" of the Faculty also recommends printing two sided)
% 11pt: font size
\documentclass[a4paper,11pt,twoside]{report}

% Uncomment this line if you want to change the default font family to palatino:
\renewcommand{\familydefault}{ppl}
% Here you can browse more fonts supported by Overleaf (just make sure to use a serif font typeface): https://www.sharelatex.com/learn/Font_typefaces

% Package needed to format table of contents (ToC):
\usepackage{tocloft}
% Indentations of sections in ToC:
\setlength{\cftsecindent}{.9cm}
% Indentations of subsections in ToC:
\setlength{\cftsubsecindent}{1.4cm}

% Uncomment the following line if you want the renderer insert space between paragraphs separated by an empty line
%\usepackage{parskip}
% ---------------------------------------------------

% Formatting chapter titles
% ---------------------------------------------------
\usepackage{titlesec}

% Declare the depth of numbering of chapters/sections/subsections:
% 0: only chapters are numbered
% 1: chapters and sections are numbered
% 2: even subsections are numbered
\setcounter{secnumdepth}{2} 

% Format text of chapter titles
% Uncommenting the line within the braces will change the chapter titles from "
%   Chapter {number}
%   {title}
% " to "
%   {number}. {title}
% "
\newcommand{\titleformatting}{
    %\titleformat{\chapter}{\normalfont\huge}{\thechapter.}{1em}{\huge\textbf}
}

% Uncommenting the following line will remove the number of the chapter before the number of the subsections (e.g. changes "2.1 {section title}" to "1. {section title}" where "2" is the number of the chapter, and "1" is the number of the section
%\renewcommand\thesection{\arabic{section}}

% Format text of section titles
%\titleformat{<which element you want to format>}[block/wrap]{<format rules for numbers of sections>}{<commands that render the numbers>}{<how much space you want between the number and the title>}{<format rules for title>}
% size rules: \large, \Large, \huge, \Huge
% more format rules: \bfseries (bold), \it (italic), \underline
% \thesection: Command to render the number of the section
\titleformat{\section}[block] {\normalfont\Large\bfseries}{\thesection.}{1em}{\Large}
% ---------------------------------------------------


% Mandatory formatting rules required by "Code of Studies and Exams" [don't change this two lines!]
% ---------------------------------------------------

% margins
\usepackage[margin=2.5cm, bindingoffset=1.25cm]{geometry}

% line spacing
\linespread{1.5}
% ---------------------------------------------------


% Useful packages
% ---------------------------------------------------

% Handle UTF8 characters
\usepackage[utf8]{inputenc}

\setlength{\emergencystretch}{10pt}

% Change numbering, e.g. to letters instead of numbers
\usepackage{enumitem}

% If you don't use \usepackage[T1]{fontenc},
%  - Words containing accented characters cannot be automatically hyphenated,
%  - You cannot properly copy-and-paste such words from the output (DVI/PS/PDF),
%  - Characters like the pipe sign, less than and greater sign give unexpected results in text.
\usepackage[T1]{fontenc}

% language specification
\usepackage[english]{babel}

% useful for equations
\usepackage{mathtools} 
% Convenience command for norm
\newcommand{\norm}[1]{\left\lVert#1\right\rVert}

\usepackage{amsmath,amssymb,amsfonts}
\DeclareMathOperator*{\argmax}{arg\,max}
\DeclareMathOperator*{\argmin}{arg\,min}
\DeclareMathOperator{\diag}{diag}

\newenvironment{tight_equations}{
    \setlength{\abovedisplayskip}{3pt}
    \setlength{\belowdisplayskip}{3pt}
}

\newcommand{\citationneeded}{~[\color{red}citation needed\color{black}]}

% greek letters in text mode (used in equation tags
\usepackage[euler]{textgreek}

% For theorems
\usepackage{amsthm}

\newtheoremstyle{def_thm}% name of the style to be used
  {20pt}% measure of space to leave above the theorem. E.g.: 3pt
  {}% measure of space to leave below the theorem. E.g.: 3pt
  {\itshape}% name of font to use in the body of the theorem
  {}% measure of space to indent
  {\bfseries}% name of head font
  {:}% punctuation between head and body
  { }% space after theorem head; " " = normal interword space
  {}% Manually specify head
\theoremstyle{def_thm}
\newtheorem*{definition}{Definition}
\newtheorem*{theorem}{Theorem}

\theoremstyle{remark}
\newtheorem*{remark}{Remark}
\newtheorem*{notation}{Notation}

\usepackage{algorithmic}
\usepackage{graphicx}
\usepackage{textcomp}
\usepackage[ruled,vlined]{algorithm2e}
%\usepackage{authblk}                     

% Format references by IEEE guidlines
\usepackage[style=ieee, backend=biber]{biblatex}
\renewbibmacro*{doi+eprint+url}{%
  \iftoggle{bbx:doi}
    {\iffieldundef{url}{\printfield{doi}}{}}
    {}%
  \newunit\newblock
  \iftoggle{bbx:eprint}
    {\usebibmacro{eprint}}
    {}%
  \newunit\newblock
  \iftoggle{bbx:url}
    {\usebibmacro{url+urldate}}
    {}}

\DeclareBibliographyAlias{artwork}{misc}

% Needed for positioning the logo of the Faculty on the first page
\usepackage[export]{adjustbox}

% Display SI units
\usepackage[binary-units=true]{siunitx}

% vulgar fractions
\usepackage{textcomp}

% smaller font size for captions
\usepackage[font=footnotesize,labelfont=bf]{caption}

% Importing the hyperref package all cross-referenced elements become hyperlinked. For example, the lines in the table of contents become links to the corresponding pages in the document.
\usepackage[hidelinks, unicode, pdfusetitle]{hyperref}

% PDF bookmarks (smart pdf viewers can show this bookmarks helping the reader to navigate inside the document)
\usepackage{bookmark}

% Support for set background colors for table cells
\usepackage[table,xcdraw]{xcolor}
% colorful text
\usepackage{xcolor}

% Transform URL-s in the text clickable automatically
\usepackage{url}

% Support for long quotes
\usepackage{csquotes}

% Support for embedded source codes
\usepackage{listings}
\usepackage{sourcecodepro} % a nice font family for code fragments
\lstset{captionpos=b, numberbychapter=false, basicstyle=\ttfamily, showstringspaces=false, columns=fullflexible}

% Tell the renderer where to look for images
\graphicspath{ {images/} }

% File holding references
\addbibresource{references.bib}
\addbibresource{manually_added_references.bib}

% Generate Dummy Text (feel free to remove it)
\usepackage{lipsum}

\usepackage{etoolbox}
%\AtBeginEnvironment{quote}{\singlespacing\small}

% ---------------------------------------------------

% First page + empty page after it
% You should bind the the original copy of the Diploma Thesis Proposal Form (collected from the Registrar’s Office) right after this blank page!
% ---------------------------------------------------

\author{\name \\ \program}
\title{\Huge{\thesisType}\\[1cm]
    \huge{\title}}
\date{\the\year}

\newcommand{\titlePage}{
    \includegraphics[valign=m, width=50pt]{ITK_logo} \parbox[c]{0.8\textwidth}{
    Pázmány Péter Catholic University\\
    Faculty of Information Technology and Bionics}
    \vspace*{\fill}
    
    {\let\newpage\relax\maketitle}
    \vspace*{\fill}
    \begin{center}
    \bigskip
    
    \supervisors
    \end{center}
}

\newcommand{\blankPage}{
    \begingroup
        \pagestyle{empty}
        \cleardoublepage
    \endgroup
    \clearpage
}
% ---------------------------------------------------

% ------------- Dokumentum legenerálása -------------
\begin{document}

\titleformatting

% This file is responsible to collect all content and assemble the document.

% Edit these lines...
\def\thesisType{Master Thesis}
\def\name{HAKKEL Tamás}
\def\program{Computer Science Engineering MSc}
\def\title{Iteratively Reweighted Algorithms for Dynamic MRI}
\def\supervisors{Supervisor: Claudio M. VERDUN \\ Faculty mentor: Dr. OLÁH András}

% ---------------------------------------------

% Render title page
\titlePage

% Set page numbering to roman numbers for all pages before the first chapter. Change it to "\pagenumbering{Roman}" if you want upper case roman numbers
\pagenumbering{roman}

% Mandatory parts
% If you don't want these mandatory parts to be included in the table of contents, simply delete or comment out the lines starting with "\addcontentsline".
\addcontentsline{toc}{chapter}{Diploma Thesis Proposal}

% If you print your thesis two sided, the you need this blank page after the Title page
\blankPage

% The "Diploma Thesis Proposal" form takes the first two page numbers, so the "Thesis Authenticity Statement" gets the page number 3.
\setcounter{page}{3}

\addcontentsline{toc}{chapter}{Thesis Authenticity Statement}
% You have nothing to do here

\chapter*{Thesis Authenticity Statement}
I, undersigned \name, student of the Faculty of Information Technology and Bionics at the Pázmány Péter Catholic University, hereby certify that this thesis was written without any unauthorized help, solely by me and I used only the referenced sources. Every part, which is quoted exactly or in a paraphrased manner, is indicated clearly with a reference. I have not submitted this thesis anywhere else.

% Signature line
\begin{flushright}
	\vspace*{.5cm}\par\noindent\makebox[2.5in]{\hrulefill}
	\par\noindent\makebox[2.5in][c]{\name}
\end{flushright}

\clearpage
\addcontentsline{toc}{chapter}{Abstract}
\chapter*{Abstract}

Although MR imagig was invented less then 50 years ago, it has revolutionized medical imaging and diagnostic process as we know it. Its versatility makes it fit a wide range of use cases, and compared to other imaging technologies, MRI demonstrates important advantages in many cases, such as lack of ionizing radiation, adjustable contrast range, and excellent differentiation between soft tissues. However, it has also a disadvantage: the imaging speed is undesirably slow. While it is only inconvenience in some cases, dynamic images of chest, for example, gets blurred by the different motions of the body, like breathing and cardiac motion. There are new technologies that improves the imaging process itself, but strong image reconstruction algorithms can vastly improve the quality.

Among the many software techniques invented to improve image quality, compressed sensing is inevitably is one of the most impactful theoretical construction, introduced by Donoho, Candès, Romberg, and Tao in 2004~\cite{candes_robust_2006, donoho_compressed_2006, candes_nearoptimal_2006}. In contrast to the Nyquist-Shannon theorem that asserts that continuous band-limited signals can be perfectly reconstructed from samples taken at a rate of twice the highest frequency present in the signal of interest, compressed sensing allows lossless reconstruction from much lower number of samples given that certain natural conditions are satisfied.

In this thesis work, we consider the classic results as well as the recent advances within of the compressed sensing framework and their application to real life MR imaging. In particular, we closely examine two recent publications presenting state-of-the-art solutions combining conventional techniques with novel ideas.

Afterwards, we present our implementation of these algorithms along with the implementation of a recently invented algorithm from the family of iteratively reweighted least squares (IRLS) methods that previously have not been applied to MRI setting yet. For the language of implementation we have chosen Julia, a new open-source language released in 2012, as this language fits well the image reconstruction problem being inherently fast with a speed often comparable to C, and providing convenient environment for fast prototyping.

Finally, we compare these algorithm with respect to reconstruction power from massively undersampled data and noise tolerance. Our results demonstrates the fitness of the Julia language to prototyping and also shows strong benchmark suggesting that a robust extension of the new IRLS method can bring enormous improvement to reconstruction quality or imaging speed.

\clearpage
% Uncomment the following line if you want to include an "Acknowledgements" page after the abstract.
%\addcontentsline{toc}{chapter}{Acknowledgements}
%\chapter*{Acknowledgements}
% Optional part to give thanks to people contributing to your work (other than your supervisor and mentor) https://seleninevcilikhayati.com/accounting-dissertation-help/

\paragraph{}
\lipsum[2] % 1 paragraph of dummy text - replace it with your thanksgiving lines
\clearpage

% Render table of contents
\tableofcontents
\clearpage

% From here page numbers are arabic numbers starting again from 1
\pagenumbering{arabic}

% Include files holding the content of your thesis -- feel free to change as you like it
\chapter{Introduction}
% The Code of Studies and Exams recommends the following content for this chapter: "Context of the work, professional justification/explanation of the relevance, motivation, objective, and a summary of the tasks performed."
% However, you are free to structure the content of your thesis as you want.

While the fast evolution of technology profoundly changed today's medicine, unarguably the medical imaging is of the fields which profited the most of the computation power recently became available. And as X-ray radiographs revolutionized medical treatments in the beginning of the 20th century,  the appearance of computer-aided imaging techniques such as computer tomography (CT), diagnostic ultrasonography, positron emisson tomography (PET), and magnetic resonance imaging (MRI) opened a new horizon drastically increasing the resolution, allowing 3D imaging, providing reliable dynamic recordings, and enhancing images by automated post-processing. In the recent decades radiology evolved to be an interdisciplinary field involving, for instance, molecular biology, nuclear physics, applied mathematics, and computer science besides the classical medical fields such as anatomy, angiology, and cardiology.

\section{Magnetic Resonance Imaging}
\subsection{Historical Background}
The beginning of the history of medical images dates back to November of 1895 when Wilhelm Conrad Röntgen discovered X-ray radiation. Besides the remarkable fact that he was awarded the first Nobel Prize in Physics in 1901, it also signifies the importance of his discovery that up until the 1960's it was the only medical imaging technique available. During the first 60 years of X-ray radiography, it underwent a remarkable development gradually increasing the resolution of images and decreasing the radiation dose that affected both the patients' and the doctors' health. Since the 1920s, visualization of motions (i.e. dynamic imaging) became possible by flouroscopy, although only to a limited extent~\cite{bradley_history_2008}.

While the theoretical background of later advances in medical images were present much earlier, the the introduction to clinical medicine took place only when the computers became powerful enough for image reconstruction tasks concerning both the computational power and the available memory. Therefore, the major breakthrough came only in the 1960-70s, when suddenly multiple imaging techniques were introduced. The first of these methods was the ultrasound imaging devised by Floyd Firestone in 1940 to detect internal flaws in metal castings~\cite{singh_origin_2007}, and used first for medical purposes in 1949 by John Wild~\cite{watts_john_2009}, but it was not until 1961 when David Robinson and George Kossoff developed the first commercially practical water path ultrasonic scanner~\cite{griffiths_historical_nodate}. Also, the concept of emission and transmission tomography was propsed in the late 1950s by David E. Kuhl, Luke Chapman and Roy Edwards, but computer tomography was invented in 1972 by Godfrey Hounsfield and Allan Cormack~\cite{richmond_sir_2004}, and the first PET camera was built for human studies by Edward Hoffman, Michael M. Ter-Pogossian, and Michael E. Phelps in 1973~\cite{noauthor_us_nodate}. Moreover, nuclear magnetic resonance (NMR) imaging was discovered by Felix Bloch and Edward Purcell in 1946, and Reymond Vahan Damadian proposed the first MR body scanner in 1969. Then Paul Lauterbur had the idea of applying magnetic field gradients in all three dimensions and a back-projection technique to create images in 1971, and he also published the first MRI images: water tubes, a living clam, and the thoracic cavity of a mouse in 1973 and 1974~\cite{rinck_short_2008}. Since then, these computer-aided methods continue to develop at a very fast pace, becoming essential tools for today's practitioners.

\subsection{Mechanism}


\subsection{Diagnostics}
\subsection{Advantages and Limitations}

\section{Compressed Sensing}
\subsection{Core Concept}
\subsection{Application Areas}
\subsection{Suitability to MRI}

\section{Julia Language}
\subsection{Objectives of The Language}
\subsection{Suitability to Our Task}

\section{Objective}

% ----------------------------------------------------
\section{Outline}
The summary of the chapters of the thesis work:

\paragraph{Chapter 2} This chapter describes something and here I summarize it in a couple sentences.

\paragraph{Chapter 3} This chapter describes something and here I summarize it in a couple sentences.

\paragraph{Chapter 4} This chapter describes something and here I summarize it in a couple sentences.

\paragraph{Chapter 5} This chapter describes something and here I summarize it in a couple sentences.

\clearpage % You need \clearpage at the end of every chapter to force images included in this chapter to be rendered in somewhere else
\chapter{Related Works}
% The Code of Studies and Exams recommends the following content for this chapter: "Literature review, short description of similar/related applications/solutions, conclusions of literature review. Independent, critical analysis of the literature."
% However, you are free to structure the content of your thesis as you want

\section{Compressed Sensing in MRI}
"Compressed Sensing MRI" (2008) by Michael Lustig, David L. Donoho, Juan M. Santos, and John M. Pauly --> "A look at how CS can improve on current imaging techniques"

"Sparse MRI: The Application of Compressed Sensing for Rapid MR Imaging" (2007) by Michael Lustig, David Donoho, and John M. Pauly

\section{Parallel Imaging}
"Self-Calibrating Nonlinear Reconstruction Algorithms for Variable Density Sampling and Parallel Reception MRI" by Loubna El Gueddari, C. Lazarus, H Carrié, A. Vignaud, Ph Ciuciu

SENSE and ESPIRiT for sensitivity map estimation

\section{First Order Gradient Methods}

My slideshow about first order methods: \url{https://docs.google.com/presentation/d/1tIZSSfzzHUgo9tlKlmhSUCCWD3ynzyZwnWtpBUe--6g/edit}

\subsection{Origins}
From Cauchy to recent explosion of optimization algorithm.

Why would second order methods speed up convergence and why aren't we able to apply them to our problem.

\subsection{Gradient Descent vs Conjugate Gradient}
How do they work, and when can conjugate gradient outperform gradient descent

\subsection{Proximal Methods and ISTA}
Proximal operator, Iterative Shrinkage/Soft Thresholding Algorithm, Forward-Backward Splitting with linear line search, Wolfe conditions

\subsection{FISTA by Nesterov}
Momentum rule and its role in speeding up convergence

\subsection{POGM by Kim and Fessler}
Why is it faster than FISTA, and why is it optimal

D. Kim and J. A. Fessler, “Optimized first-order methods for smooth convex minimization,” Math. Program., vol. 159, no. 1, pp. 81–107, Sep. 2016, doi: 10.1007/s10107-015-0949-3.

\subsection{ADMM}
what is it and why is it good

\section{Decompositions}

\subsection{Low rank and Sparse}
C. Y. Lin and J. A. Fessler, “Efficient Dynamic Parallel MRI Reconstruction for the Low-Rank Plus Sparse Model,” IEEE Transactions on Computational Imaging, vol. 5, no. 1, pp. 17–26, Mar. 2019, doi: 10.1109/TCI.2018.2882089.

J. A. Fessler, “Optimization methods for MR image reconstruction (long version),” arXiv:1903.03510 [eess, math], Jun. 2019.

\subsection{Multiscale}

Ong's dissertation: “Low Dimensional Methods for High Dimensional Magnetic Resonance Imaging,” 2018.

F. Ong et al., “Extreme MRI: Large-Scale Volumetric Dynamic Imaging from Continuous Non-Gated Acquisitions,” arXiv:1909.13482 [physics], Dec. 2019.

Differences between Ong's dissertation and his "extreme MRI" preprint paper

\section{IRSL}
Iteratively reweighted Least Squares method

C. Kümmerle and C. M. Verdun, “Denoising and Completion of Structured Low-Rank Matrices via Iteratively Reweighted Least Squares,” arXiv:1811.07472 [cs, math], Nov. 2018.

Henry Adams, Lara Kassab, and Deanna Needell "An Iterative Method for Structured Matrix Completion"

\clearpage % You need \clearpage at the end of every chapter to force images included in this chapter to be rendered in somewhere else
\chapter{Methods}
% The Code of Studies and Exams recommends the following content for this chapter: "Presentation of used methodologies/technologies: In line with the topic of the dissertation, the professional background related to the solution and implementation of the task must be detailed."
% However, you are free to structure the content of your thesis as you want

\section{FunctionOperators package}
Describe this package what I created, and explain why was it useful.
\subsection{Objective}
\subsection{Features}

\section{Implementation of Sparse+Low Rank algorithms}
Re-implementation of Lin \& Fessler paper
\subsection{PINCAT dataset}
\subsection{Description of Algorithms in the article}

\section{Implementation of Multiscale Decomposition}

\subsection{NUFFT}
\paragraph{What is NUFFT}
\paragraph{Available solutions}
\paragraph{My implementation}

\subsection{Description of algorithm}

\subsection{Optimization possibilities}
\begin{enumerate}
    \item Batch-processing: compute NUFFT of all channels at once
    \item Parallelization:
    \begin{enumerate}
        \item NUFFT
        \item Algorithm
    \end{enumerate}
    \item GPU-specific optimizations:
    \begin{enumerate}
        \item Blocking operator might cause CPU bottleneck
    \end{enumerate}
\end{enumerate}

\subsection{Parallelization Solution}
Details how I made the code run parallel

\section{IRLS Implementation}

\subsection{Description of Algorithm}
\subsection{Implementation Details}

\clearpage % You need \clearpage at the end of every chapter to force images included in this chapter to be rendered in somewhere else

\chapter{Results}
% The Code of Studies and Exams recommends the following content for this chapter: "Evaluation, critical analysis of the implemented technical solutions, possibilities for further development."
% However, you are free to structure the content of your thesis as you want

\section{Sparse + Low Rank Algorithm}
\subsection{Running Speed and Memory Used}
\subsection{Readability}

\section{Multi-scale Algorithm}
\subsection{Running Speed and Memory Used}

\section{ILRS}
\subsection{Comparism with Sparse+Low Rank}
\paragraph{Running speed and memory requirement}
\paragraph{Convergence speed}
\paragraph{Robustness:} Effect of error on the input data
\paragraph{Recovery capability:} Can we reach the same accuracy on recovered image using less data?

\subsection{Comparism with Multi-scale}
\paragraph{Running speed and memory requirement}
\paragraph{Convergence speed}
\paragraph{Robustness:} Effect of error on the input data
\paragraph{Recovery capability:} Can we reach the same accuracy on recovered image using less data?

\clearpage % You need \clearpage at the end of every chapter to force images included in this chapter to be rendered in somewhere else
\chapter{Summary}
% The Code of Studies and Exams recommends the following content for this chapter:
% A summary of the problems solved compared to the objectives presented in the introduction and  in the Thesis Proposal Form. Opportunities to move forward, questions motivating the future works, outlook.
% However, you are free to structure the content of your thesis as you want

\section{Objectives}
\section{Achievements}
\section{Future Plans}

% Bibliography
\printbibliography
\addcontentsline{toc}{chapter}{Bibliography}

% Appendices -- if you don't have any, just delete the following two lines
%\appendix
%\chapter{Appendix}

Here you can present all the materials that helps the understanding of your work, but either 1) not your work, 2) not necessary to understanding, 3) simply just interesting things not directly related to your topic, or 4) your text already exceeded 120 pages, and need to make it shorter by moving parts to appendix... :D Anyway, if you have a large amount of images, code, or measurement data, it is required to insert them here rather than in chapters. You can have multiple appendices (possible separated into multiple files) according to the content to be attached.


\end{document}

